\section*{Chapitre \ref{ch:b3:complex-mechanisms} --- Mécanismes réactionnels des complexes}

\begin{solution}{exo:b3:complex-mechanisms:1}
Inertes~: $\ce{[Cr(H2O)6]^{3+}}$ ($d^3$) et $\ce{[Co(NH3)6]^{3+}}$ ($d^6$ bas spin), grandes énergies
d'activation de champ des ligands. \omterm{def:b3:complex-mechanisms:labile}{Labiles}~: $\ce{[Cr(H2O)6]^{2+}}$ ($d^4$ haut spin) et $\ce{[Cu(H2O)6]^{2+}}$ ($d^9$),
déformés par l'effet Jahn--Teller, avec de longues liaisons axiales
faibles~; $\ce{[Zn(H2O)6]^{2+}}$ ($d^{10}$), sans barrière de champ des
ligands.
\end{solution}

\begin{solution}{exo:b3:complex-mechanisms:2}
Lorsque $k_2[\mathrm Y] \gg k_{-1}[\mathrm X]$, le dénominateur vaut $k_2[\mathrm Y]$ et $v = k_1[\mathrm{ML_5X}]$~; lorsque $k_{-1}[\mathrm X] \gg k_2[\mathrm Y]$, il vaut
$k_{-1}[\mathrm X]$ et la seconde forme s'ensuit (un pré-équilibre inhibé
par le ligand partant).
\end{solution}

\begin{solution}{exo:b3:complex-mechanisms:3}
D~: $\Delta V^{\ddagger} > 0$ et $\Delta^{\ddagger}S^\circ > 0$ (un ligand part). A~: tous deux négatifs (un
ligand est lié dans l'état de transition).
\end{solution}

\begin{solution}{exo:b3:complex-mechanisms:4}
$cis$-$\ce{[PtCl2(NH3)2]}$ à partir de $\ce{[PtCl4]^{2-}}$~; $trans$-$\ce{[PtCl2(NH3)2]}$ à partir de $\ce{[Pt(NH3)4]^{2+}}$
(l'\omterm{def:b3:complex-mechanisms:trans-effect}{effet trans} de \ce{Cl-} dépasse celui de \ce{NH3}).
\end{solution}

\begin{solution}{exo:b3:complex-mechanisms:5}
$k_{\mathrm{obs}} = 2.0 \times \num{3.0e4} \times 0.10/1.2 = \qty{5.0e3}{s^{-1}}$~; à \qty{1.0}{mol/L}, $\num{6.0e4}/3.0 = \qty{2.0e4}{s^{-1}}$~; limite
$k_i = \qty{3.0e4}{s^{-1}}$.
\end{solution}

\begin{solution}{exo:b3:complex-mechanisms:6}
$\Delta V^{\ddagger} = -RT\ln2/\Delta p = -8.314 \times 298 \times 0.693/\num{99.9e6} = \qty{-1.7e-5}{m^3.mol^{-1}} = \qty{-17}{cm^3.mol^{-1}}$~:
\omterm{def:b3:complex-mechanisms:mechanisms}{mécanisme associatif}.
\end{solution}

\begin{solution}{exo:b3:complex-mechanisms:7}
La phosphine, fort donneur $\sigma$ et accepteur $\pi$, affaiblit la liaison en trans (\omterm{def:b3:complex-mechanisms:trans-effect}{influence
trans}, liaison plus longue) et la labilise (\omterm{def:b3:complex-mechanisms:trans-effect}{effet trans})~: ce chlorure est
remplacé en premier.
\end{solution}

\begin{solution}{exo:b3:complex-mechanisms:8}
Le ligand pontant se retrouve transféré au métal oxydé~; la vitesse dépend
fortement de la nature du pont potentiel~; un intermédiaire binucléaire est
détecté~; et la vitesse ne peut pas dépasser celle de la substitution
nécessaire pour former le pont sur le partenaire \omterm{def:b3:complex-mechanisms:labile}{labile}.
\end{solution}

\begin{solution}{exo:b3:complex-mechanisms:9}
$(\lambda + \Delta G^\circ)^2/4\lambda$ avec $4\lambda = 320$~: 20.0, 11.3, 0 et \qty{5.0}{kJ/mol} (la dernière dans la
\omterm{def:b3:complex-mechanisms:inverted}{région inversée}).
\end{solution}

\begin{solution}{exo:b3:complex-mechanisms:10}
$k_{12} = \sqrt{4.0 \times \num{2.0e5} \times \num{1.0e4}} = \qty{8.9e4}{L.mol^{-1}.s^{-1}}$.
\end{solution}

\begin{solution}{exo:b3:complex-mechanisms:11}
$d^5$ haut spin~: 0~; $d^7$~: $1.33\,Dq$~; $d^8$~: $2.0\,Dq$. Vitesses d'échange d'eau attendues~: \ce{Mn^{2+}} $>$ \ce{Co^{2+}} $>$
\ce{Ni^{2+}}, l'ordre observé.
\end{solution}

\begin{solution}{exo:b3:complex-mechanisms:12}
Pour des réactions très exergoniques, la vitesse intrinsèque dépasse la
vitesse de rencontre des ions~: la constante observée reste à la limite de
diffusion et la diminution est masquée. Avec le donneur et l'accepteur
fixés à une distance donnée dans une même molécule, il n'y a pas d'étape de
diffusion, et une série d'accepteurs de force motrice croissante a montré
la vitesse qui augmente, passe par un maximum et diminue.
\end{solution}

\begin{solution}{pb:b3:complex-mechanisms:1}
\textbf{1.} $5d^8$.
\textbf{2.} Le champ fort d'un ion $5d$ rend très haute l'orbitale $d_{x^2-y^2}$, qui pointe vers les quatre
ligands~: les huit électrons s'apparient dans les quatre orbitales plus basses.
\textbf{3.} Le complexe a 16 électrons de valence et un site axial libre~:
le ligand entrant se lie d'abord, via un état de transition pentacoordiné.
\textbf{4.} $k_{\mathrm{obs}} = k_1 + k_2[\mathrm Y]$~: un chemin par le solvant (attaque du solvant, puis remplacement rapide par Y)
et un chemin direct.
\textbf{5.} La substitution prend des heures (partie III)~: inerte à
l'échelle de la minute, pas à celle de la journée.
\textbf{6.} Le cisplatine, car la deuxième molécule de \ce{NH3} remplace un chlorure en trans d'un chlorure
(\omterm{def:b3:complex-mechanisms:trans-effect}{effet trans} $\ce{Cl-} > \ce{NH3}$)~; il se forme aussi des sous-produits.
\textbf{7.} L'iodure a un \omterm{def:b3:complex-mechanisms:trans-effect}{effet trans} plus fort que le chlorure, ce qui rend
l'effet orienteur plus net et plus rapide.
\textbf{8.} Dans $\ce{[PtI3(NH3)]-}$, c'est l'iodure en trans d'un iodure qui est labilisé, et non celui en trans de
\ce{NH3}~: la deuxième molécule d'ammoniac entre en cis de la première.
\textbf{9.} L'iodure d'argent précipite et $cis$-$\ce{[Pt(H2O)2(NH3)2]^{2+}}$ reste en solution.
\textbf{10.} Le chlorure remplace les deux molécules d'eau à leurs propres
positions~; les ligands ammine ne sont pas touchés.
\textbf{11.} À partir de $\ce{[Pt(NH3)4]^{2+}}$ et de deux chlorures.
\textbf{12.} Ses deux positions réactives sont opposées~: il ne peut pas
lier deux bases voisines d'un même brin d'ADN, la lésion par laquelle agit
le cisplatine.
\textbf{13.} $\ln2/\num{9.0e-5} = \qty{7.7e3}{s} = \qty{2.1}{h}$.
\textbf{14.} $1 - \eu^{-\num{9.0e-5} \times 7200} = 0.48$.
\textbf{15.} L'eau est un bien meilleur groupe partant que le chlorure~:
l'aquacomplexe est rapidement substitué par l'azote d'une base de l'ADN.
\textbf{16.} Associative, comme toutes les substitutions du platine(II)~:
\omterm{def:b3:complex-mechanisms:activation-volume}{volume d'activation} négatif.
\textbf{17.} La forte concentration en chlorure ramène l'équilibre vers la
forme dichlorée et garde le médicament intact jusqu'à ce qu'il atteigne les
cellules.
\textbf{18.} $f = K/(K + [\ce{Cl-}])$.
\textbf{19.} $\num{3.0e-3}/0.103 = 0.029$.
\textbf{20.} $\num{3.0e-3}/0.0070 = 0.43$.
\textbf{21.} 15.
\textbf{22.} Ce n'est que dans les cellules qu'une fraction notable existe
sous la forme aqua réactive.
\textbf{23.} Les ligands soufrés, comme le glutathion et la cystéine et la
méthionine des protéines, qui lient fortement le platine, mou.
\textbf{24.} L'ammoniac se lie fortement au platine et se trouve en trans
du chlorure, un ligand de faible \omterm{def:b3:complex-mechanisms:trans-effect}{effet trans}~: les liaisons platine--ammine
ne sont pas labilisées.
\textbf{25.} \textbf{Environ 15 fois plus de médicament se trouve sous la
forme aqua réactive dans une cellule que dans le plasma.}
\end{solution}
